\documentclass{beamer}

% --------------------
% Theme and Packages
% --------------------
\usetheme{Madrid}
\usecolortheme{default}
\usepackage{graphicx}
\usepackage{booktabs}
\usepackage{amsmath}
\usepackage{hyperref}

% --------------------
% Title Information
% --------------------
\title[Kiswahili Audio Intelligence Pipeline]
{An End-to-End Kiswahili Audio Processing Pipeline for Speech Recognition, Sentiment Analysis, and Summarization}

\author{Kevin Obote}
\institute{
Strathmore University \\ iLabAfrica \\
\vspace{0.2cm}
\textbf{Supervisor:} Dr. Betsy Muriithi
}
\date{\today}

% --------------------
% Begin Document
% --------------------
\begin{document}

% --------------------
% Title Slide
% --------------------
\begin{frame}
\titlepage
\end{frame}

% --------------------
% Outline
% --------------------
\begin{frame}{Presentation Outline}
\tableofcontents
\end{frame}

% ==========================================================
% CHAPTER 1: INTRODUCTION
% ==========================================================
\section{Chapter 1: Introduction}

\begin{frame}{Background of the Study}
Kiswahili is one of the most widely spoken African languages, yet remains underrepresented in speech and language technologies.
\begin{itemize}
    \item Growth in audio-based digital content
    \item Dominance of high-resource languages in ASR systems
    \item Limited end-to-end audio intelligence systems for Kiswahili
\end{itemize}
\end{frame}

\begin{frame}{Problem Statement}
Despite advances in deep learning:
\begin{itemize}
    \item Kiswahili ASR systems exhibit high word error rates
    \item Sentiment models perform poorly on neutral and factual speech
    \item Summarization lacks reliable ground truth datasets
\end{itemize}
This limits the usability of Kiswahili audio analytics in real-world applications.
\end{frame}

\begin{frame}{Research Objectives}
\textbf{Main Objective:}
\begin{itemize}
    \item To design and deploy an end-to-end Kiswahili audio processing pipeline
\end{itemize}

\textbf{Specific Objectives:}
\begin{itemize}
    \item Develop a Kiswahili ASR system
    \item Perform sentiment analysis on speech transcripts
    \item Generate abstractive summaries from spoken content
    \item Deploy the system as a functional prototype
\end{itemize}
\end{frame}

\begin{frame}{Scope and Significance}
\begin{itemize}
    \item Focus on Kiswahili speech data
    \item Offline batch audio processing
    \item Contribution to low-resource language NLP research
\end{itemize}
\end{frame}

% ==========================================================
% CHAPTER 2: LITERATURE REVIEW
% ==========================================================
\section{Chapter 2: Literature Review}

\begin{frame}{Automatic Speech Recognition for Low-Resource Languages}
\begin{itemize}
    \item HMM-GMM vs Transformer-based ASR
    \item Wav2Vec2 and self-supervised learning
    \item Challenges in African language datasets
\end{itemize}
\end{frame}

\begin{frame}{Sentiment Analysis and Speech}
\begin{itemize}
    \item Text-based sentiment dominance
    \item Multilingual transformers
    \item Limitations in neutral sentiment detection
\end{itemize}
\end{frame}

\begin{frame}{Abstractive Text Summarization}
\begin{itemize}
    \item Extractive vs Abstractive methods
    \item T5 and mT5 architectures
    \item Evaluation challenges using ROUGE
\end{itemize}
\end{frame}

% ==========================================================
% CHAPTER 3: METHODOLOGY
% ==========================================================
\section{Chapter 3: Methodology}

\begin{frame}{Research Design}
The study follows the CRISP-DM framework:
\begin{enumerate}
    \item Business Understanding
    \item Data Understanding
    \item Data Preparation
    \item Modeling
    \item Evaluation
    \item Deployment
\end{enumerate}
\end{frame}

\begin{frame}{System Architecture}
\begin{figure}
    \centering
    \includegraphics[width=0.9\textwidth]{figures/system_architecture.pdf}
    \caption{Overall System Architecture}
\end{figure}
\end{frame}

\begin{frame}{End-to-End Processing Flow}
\begin{figure}
    \centering
    \includegraphics[width=0.9\textwidth]{figures/processing_flow_diagram.pdf}
    \caption{Audio Processing Flow Architecture}
\end{figure}
\end{frame}

\begin{frame}{Data Sources and Preparation}
\begin{itemize}
    \item Mozilla Common Voice (Kiswahili subset)
    \item Audio normalization and resampling
    \item Text cleaning and chunking
\end{itemize}
\end{frame}

% ==========================================================
% CHAPTER 4: SYSTEM IMPLEMENTATION
% ==========================================================
\section{Chapter 4: System Implementation}

\begin{frame}{ASR Model Implementation}
\begin{itemize}
    \item Wav2Vec2 fine-tuned for Kiswahili
    \item Whisper used as a baseline
    \item GPU and CPU optimized training
\end{itemize}
\end{frame}

\begin{frame}{Sentiment Analysis Implementation}
\begin{itemize}
    \item DistilBERT multilingual model
    \item Labels: Positive, Neutral, Negative
    \item Chunk-level sentiment aggregation
\end{itemize}
\end{frame}

\begin{frame}{Summarization Implementation}
\begin{itemize}
    \item T5-based abstractive summarization
    \item Text chunking for long transcripts
    \item Sentence-level summary merging
\end{itemize}
\end{frame}

% ==========================================================
% CHAPTER 5: RESULTS AND EVALUATION
% ==========================================================
\section{Chapter 5: Results and Evaluation}

\begin{frame}{ASR Evaluation Metrics}
\textbf{Word Error Rate (WER):}
\[
WER = \frac{S + D + I}{N}
\]
\end{frame}

\begin{frame}{ASR Results}
\begin{table}
\centering
\begin{tabular}{lc}
\toprule
Metric & Value \\
\midrule
Final Test WER & 0.2917 \\
Best Validation WER & 0.1068 \\
Average Sample WER & 0.1914 \\
\bottomrule
\end{tabular}
\caption{ASR Performance Results}
\end{table}
\end{frame}

\begin{frame}{Sentiment Analysis Results}
\begin{table}
\centering
\begin{tabular}{lc}
\toprule
Metric & Score \\
\midrule
Precision & 0.015 \\
Recall & 0.014 \\
F1-score & 0.0144 \\
\bottomrule
\end{tabular}
\caption{Sentiment Analysis Performance}
\end{table}
\end{frame}

\begin{frame}{Summarization Results}
\begin{table}
\centering
\begin{tabular}{lccc}
\toprule
Model & ROUGE-1 & ROUGE-2 & ROUGE-L \\
\midrule
AfriTeVa V2 & 0.300 & 0.300 & 0.300 \\
mT5 XLSum & \textbf{0.460} & \textbf{0.321} & \textbf{0.424} \\
mt5-small & 0.300 & 0.300 & 0.300 \\
\bottomrule
\end{tabular}
\caption{Summarization Evaluation Results}
\end{table}
\end{frame}

% ==========================================================
% CHAPTER 6: CONCLUSION AND FUTURE WORK
% ==========================================================
\section{Chapter 6: Conclusion and Future Work}

\begin{frame}{Conclusion}
\begin{itemize}
    \item Successfully developed an end-to-end Kiswahili audio pipeline
    \item Demonstrated feasibility of low-resource speech intelligence
    \item Deployment validates real-world applicability
\end{itemize}
\end{frame}

\begin{frame}{Limitations}
\begin{itemize}
    \item Limited labeled sentiment data
    \item Evaluation constraints for summarization
    \item High computational requirements
\end{itemize}
\end{frame}

\begin{frame}{Future Work}
\begin{itemize}
    \item Domain-specific fine-tuning
    \item Improved evaluation methodologies
    \item Mobile and offline deployment
\end{itemize}
\end{frame}

% --------------------
% Thank You
% --------------------
\begin{frame}
\centering
\Huge{Thank You}\\
\vspace{0.5cm}
\Large{Questions?}
\end{frame}

\end{document}
